\begin{frame}{Introduction}

\textbf{Virtual sensing }
\begin{itemize}	
	\justifying
	\item VS techniques, also called soft sensing, proxy sensing, inferential sensing, or surrogate sensing, are used to provide feasible and economical alternatives to costly or impractical physical measurement instrument.
	\item This technique predicts parameters using mathematical models, which are not measured using physical sensors.
	\item A VS system uses information available from other measurements and process parameters to calculate an estimate of the quantity of interest.
	\item In this work, expensive DO sensor is replaced with solubility equations. 
	\item Solubility equations used to estimate the DO: 
	\begin{itemize}
	\item The Weiss Equations \cite{WEISS1970}
	\item The Benson and Krause Equations \cite{Benson1980,Benson1984}
	\item The Garcia and Gordon Equations \cite{Garcia1992}
	\end{itemize}
\end{itemize}
\end{frame}	

\input{1.Parts/Weiss_EQ.tex}

\begin{frame}%[allowframebreaks]
\frametitle{Node S} % and working}

\begin{itemize}	
	\justifying
	\item	Implemented virtual sensing in Node S of WQM system.
	\item 	Here costly DO sensors are replaced by virtual sensing.
	\item  	Virtual sensing uses solubility equations and interdependency equations.  %(Weiss equation).
	%	\item   We also virtual sense salinity using the values from conductivity sensor.
\end{itemize}

\begin{figure}[h!]
	\includegraphics[trim={.1cm .1cm .1cm .1cm},clip,width=8cm]{Im/PPT/NodeS.pdf}
	\centering
	\caption{Block Diagram of Node S.}
	\label{Fig:NodeS_Block}
\end{figure}	
\end{frame}

\begin{frame}{Working of Node S}%[allowframebreaks]
\begin{columns}
	
	\begin{column}{0.4\textwidth}
		\begin{figure}[]
			\includegraphics[trim={.1cm .1cm .1cm .1cm},clip,width=4.25cm]{Im/PPT/fig4.pdf}
			\centering
			\caption{Working of Node S.}
			\label{Fig:NodeS_working}
		\end{figure}
	\end{column}
	
	\begin{column}{0.7\textwidth}
		\begin{itemize}
			\justifying
			\item VS is implemented in the working of Node S.
			\item First, the reading from the physical sensors \\is taken.
			\item DO is estimated using the reading from physical \\ sensors. 
		\end{itemize}
	\end{column}
	
\end{columns}		
	
\end{frame}



